\cotefile{157-5}
\scanpage{543}
\mspage{\deleted{4\illegible{1}}454}

soit cofinal dans $J$, donc dans $I$. Comme cet ens \inserted{$J'$} est contenu dans $I_{i',v}$, cela implique que $I_{i',v}$ est cofinal dans $I$, i.e. $v\in N_{i'}$. On a d'autre part, pour $j\in\deleted{\unsure{I_{i,u}|I}}J'$
\[
\underset{v}{\lfloor\!\underline{\alpha_{ji'}}\!\rfloor}\,\alpha_{i'i}=\underset{u}{\lfloor\!\underline{\alpha_{ji}}\!\rfloor}
\]
d'où $u=v\alpha_{i'i}$ avec $v\in N_{i'}$, d'où $u\in N_{\unsure{i'|i}}\alpha_{i'i}$, qed. $\big]$

\deleted{Je dis que $\unsure{R|X}$} \inserted{Pour $i\in I$, soit}
\[
\alpha_i\colon\underbrace{\mathrm{Hom}(\deleted{\unsure{x|a}}\xi\deleted{\unsure{\,\!,|i}},\underset{\substack{\Vert\\\xi}}{x_i})}_{=M}\longrightarrow
\underset{\substack{\Vert\\\underset{i}{\varinjlim}\,\mathrm{Hom}(\xi,x_i)}}{\underline{\mathrm{Hom}(\deleted{\unsure{x|a}}\xi,F)=F(\xi)}}
\]
le morphisme canonique.

On trouve ceci :

\noindent\makebox[0pt][r]{(4.15.8.4)\quad}Si $u\in N_i$, alors la relation d'équivalence $R_i$ dans $M$ définie par $\alpha_i$ ($v\overset{R_i}{\sim}v'\iff\alpha_i(\unsure{v|w})=\alpha_i(v')$) \deleted{\unsure{qui s'égale|équivaut}} \inserted{est égale} à celle définie par l'application $v\longmapsto uv$ de $M$ dans $M$, i.e.
\[
\alpha_i(v)=\alpha_i(v')\quad\Longleftrightarrow\quad uv=uv'.
\]

C'est immédiat : puisque $I_{i,u}$ \unsure{est} cofinal dans $I$, \deleted{\unsure{l'|il}} on a $\alpha_i(v)=\alpha_i(v')$ sss $\exists\, j\in I_{i,u}$ avec $\alpha_{ji}v=\alpha_{j\unsure{i|{}}}v'$, or $\alpha_{ji}=u$.

\deleted{\unsure{Notons|Posons} $R_u$}

Pour $u\in M$, notons $R_u$ la relation d'équivalence dans $M$ définie par $v\longmapsto uv$. L'ens des relations d'équivalence dans
% archive sheet number: 542
\scanpage{544}
\mspage{455}

\unsure{obtenu|obtenues} ainsi est de cardinal $\leq\operatorname{card} M\leq\pi_0$, donc l'ens des relations d'équivalence $R_i$ $(i\in I)$ est lui aussi de cardinal $\leq\pi_0$. Cela montre, $\unsure{\pi|\pi_0}$ étant \unsure{grd} devant $\pi_0$, qu'il existe une relation d'équivalence \inserted{de la forme} $R_u$ $(u\in M)$, et un ens. cofinal $I'$ de $I$, tel que $\forall i\in I'$, on ait $R_i=R_u$. Cela montre que, quitte
à remplacer $I$ par l'ens ordonné cofinal $I'$, qu'on peut supposer que l'on a
\[
R_i=R_u\qquad\forall i\in I.
\tag*{\makebox[0pt][r]{\makebox[\displaywidth][l]{(4.15.8.5)}}}
\]
Soit alors
\[
\deleted{\unsure{N|V}}\unsure{E|N}=M/R_u=M/R_i\qquad(\forall i\in I),
\]
on voit donc que les $\alpha_{ji}\in M$ \inserted{(pour $i\leq j$)} sont tels que l'application
\[
\tau_{\alpha_{ji}}\colon v\longmapsto\alpha_{ji}v:M\mathbin{\unsure{-|\to}}M
\]
est compatible avec la relation d'équivalence $R_u$, donc passe au quotient en une
application $\overline{\alpha_{ji}}\colon\deleted{N}E\longrightarrow\deleted{N}E$, \deleted{\unsure{et}} laquelle est d'ailleurs \uline{injective}.

Soit $M'$ le sous-monoïde de $M$ formé des $u\in M$ qui ont la propriété \inserted{de stabilité} précédente, \deleted{et $\overline{\deleted{\unsure{N|E}}M'}$ le sous-monoïde} \inserted{ainsi que la propriété d'injectivité} qu'on vient de dire, et $\overline{M'}$ le monoïde \deleted{\unsure{quoti}} image de $M'$ dans le monoïde
\[
\operatorname{Mon}(\deleted{N}E)\subset\operatorname{End}_{\unsure{\mathrm{Ens}|\mathrm{En}}}(\deleted{N}E)
\]
% archive sheet number: 543
\scanpage{545}
\mspage{456}

des applications injectives de $\deleted{\unsure{N|V}}\unsure{E|N}$ dans lui-même. On voit donc que le système inductif donné, défini par les $\alpha_{ji}$, prend ses valeurs dans le sous-monoïde $M'$ de $M$ (i.e. dans la sous-catégorie
correspondante de \deleted{\unsure{B}} la catégorie $B_M$), et que la \deleted{système} \inserted{limite} inducti\deleted{f}\inserted{ve} $F(\xi)$ qui nous intéresse (pour prouver $\operatorname{card} F(\xi)\leq\pi_0$)
est \deleted{\unsure{une}} \inserted{isomorphe à} la limite inductive du système inductif \deleted{\inserted{\unsure{de|dans} la catégorie} de \unsure{type|défini par|défini}} d'un \deleted{\unsure{défini}} \inserted{\unsure{système|répétition}} d'ens tous égaux à $\deleted{N}E$, les morphismes de transition étant les $\overline{\alpha_{ji}}\in\overline{M'}\subset\operatorname{Mon}(\deleted{N}E)$. De plus,
\deleted{\unsure{comme} \inserted{\unsure{on}}} \deleted{$\unsure{M|M'|E}$} \unsure{on a}
\[
\operatorname{card}\overline{M'}\leq\operatorname{card}M\leq\pi_0.
\]
\deleted{Ainsi,} Le résultat \inserted{$\operatorname{card}F(\xi)\leq\pi_0$} voulu résultera alors du

\uline{Lemme}\inserted{-clef} 4.15.8.1. Soit $M$ un monoïde de monomorphismes (i.e tel que $uv=uv'\Longrightarrow v=v'$, pour $u,v,v'\in M$), et soit
$(\alpha_{ji})_{i\leq j}$ un système inductif dans $B_M$, indexé par un ens. ordonné $I$ \unsure{grd} devant $\pi_0$, où $\pi_0$ est un cardinal
% archive sheet number: 544
\scanpage{546}
\mspage{457}

infini tel que $\operatorname{card} M\leq\pi_0$. Alors il existe un $i_0\in I$ tel que pour \unsure{tt} $i\geq i_0$, $\alpha_{i_0,i}$ soit un iso.

Si on applique ce lemme au cas de $\overline{M'}$ et des $\overline{\alpha_{ji}}$, on trouve donc qu'il existe un $i_0\in I$ tel que les $\overline{\alpha_{ji_0}}$ soient des iso. Cela implique que
$F(\xi)\simeq\underset{i}{\varinjlim}\,\underset{E}{\lfloor\!\underline{\deleted{\unsure{N|V}}\unsure{E|N}_i}\!\rfloor}$ est isomorphe à $E_{i_0}$, donc est de cardinal égal à $\operatorname{card} \unsure{E_{i_0}|E)_{i_0}}=\operatorname{card} E\leq\pi_0$, ce qui achève la démonstration.

D'ailleurs, \unsure{l'|le'}\deleted{\unsure{idée|id}}énoncé \unsure{4.8.13|4.5.13} implique \unsure{4.15.8.1|4.16.8.1}, car dans le système inductif des $(M_i=M,\alpha_{ji})$, de $\varinjlim$ égale à $F(\xi)$, \unsure{S'il|s'il} n'existait pas de $i_0$ comme stipulé dans
4.15.8.1, alors on en conclurait \inserted{aisément} que l'on aurait $\operatorname{Card} F(\xi)>\pi_0$, contrairement à \unsure{4.8.13|4.18.13}. Nous allons dégager ceci dans un lemme :
% archive sheet number: 545
\scanpage{547}
\mspage{458}

\uline{Lemme} 4.1\deleted{\unsure{6}}5.8.2. Soit $(E_i)_{i\in\unsure{J|I}}$ un système inductif d'ensembles, indexé par un ens. ordonné \inserted{$I$} grand devant $\pi_0$. On suppose qu'il existe $i_0\in I$ tel que

a) $i\geq i_0$ implique $\operatorname{card} E_i=\pi_0$, et

b) $j\geq i\geq i_0$ implique que $\deleted{\unsure{\alpha}}\unsure{u|\alpha}_{ji}:E_i\mathbin{\unsure{-|\to}} E_j$ est un mono. Supposons de plus

c) $\forall\,\deleted{\unsure{j_0|j}}i\deleted{\unsure{\in|>}}\geq \unsure{I_0|i_0}$, $\exists\, j>\deleted{\unsure{j_0|j}}i$ tel que $\deleted{\unsure{\alpha}}u_{j\deleted{\unsure{j_0|i_0}}i}$ ne soit \uline{pas} iso. (i.e. ne soit pas surjectif).

\deleted{Alors} Soit $E_\infty=\underset{i\in I}{\unsure{\underline{\lim}|\varinjlim}}\,E_i$. Alors
\[
\operatorname{card} E_\infty>\pi_0.
\]

Ce lemme implique donc l'\uline{équivalence} de 4.8.13 et du ``lemme-clef'' 4.15.8.1.

Pour prouver 4.15.8.2, nous allons prouver quelque chose d'un peu plus précis :

\uline{Corollaire} 4.15.8.3. Sous les conditions de \unsure{4.15.8.2|4.35.8.2}, soit $\alpha_0$ le premier ordinal tel que $\operatorname{card}\alpha_0>\pi_0$, et soit $I_{\alpha_0}$ l'ens des ordinaux $\beta<\alpha$ (ensemble \inserted{ordonné} qui est grand devant $\pi_0$). Il existe alors une application

\inserted{\unsure{str.}} croissante $I_{\alpha_0}\longrightarrow I$, $\alpha\longmapsto i_\alpha$, \inserted{appliquant $0$ dans $i_0$}
% archive sheet number: 546
\scanpage{548}
\mspage{459}

\deleted{\unsure{tel|telle}} que \deleted{l'on ait} \inserted{$i_0$,} pour \unsure{tt} $\alpha\in I_{\alpha_0}$, on ait
\[
\underset{\substack{\text{c'est un \uline{mono}}\\\text{d'après b, vu que}\\i_{\alpha+1}>i_\alpha\geq i_0}}{\underline{u_{i_{\alpha+1},i_\alpha}}}\quad\text{n'est pas iso.}
\]
\margin{\makebox[0pt][r]{\smash{\rule[-3\baselineskip]{0.4pt}{3\baselineskip}}\quad}\uline{mieux} : $\forall\alpha\in I_{\alpha_0}$, $\deleted{\unsure{E}}\underset{\beta\in I_\alpha}{\varinjlim}\,E_{i_\beta}\to E_{i_\alpha}$ n'est pas surjectif}

On a alors
\[
E'_\infty\deleted{\unsure{E_\infty}}\overset{\mathrm{def}}{=}\underbrace{\underset{\alpha\in I_{\alpha_0}}{\varinjlim}\,E_{i_\alpha}}\lhook\joinrel\longrightarrow
\underset{i\in I}{\lim}\,E_i\overset{\mathrm{def}}{=}E_\infty
\]
et \deleted{$\unsure{E_\alpha}$}.
\[
\operatorname{Card} E'_\alpha\geq\underset{\substack{\Vert\\\text{successeur}\\\text{de }\pi_0}}{\underline{\operatorname{Card}\,\deleted{\unsure{\alpha}}I_{\alpha_0}}}>\pi_0,
\]
d'où $\operatorname{card} E_\infty\geq\operatorname{card} E'_\infty>\pi_0$.

\inserted{\uline{Dém}} La construction des $i_\alpha$ par récurrence transfinie est immédiate, en utilisant c) et le fait que $I$ est \unsure{grd|grand} devant $\pi_0$, et que pour tt $\alpha\in I_{\alpha_0}$, l'ens
$I_\alpha$ des \deleted{\unsure{ens|ord|$\alpha$}} $\beta<\alpha$ est de cardinal $\leq\pi_0$. Le fait que $E'_\infty\mathbin{\unsure{-|\to}}E_\infty$ soit injectif résulte de b). Le fait que \deleted{$\operatorname{card} E$}
\[
\operatorname{card} E'_\infty\geq\operatorname{card} I_{\alpha_0}
\]
résulte du fait que $\deleted{\unsure{E_\alpha|E_\infty}}E'_\infty$ est réunion croissante de sous-ensembles identifiés à $E_{i_\alpha}$, et pour tt
% archive sheet number: 547
\scanpage{549}
\mspage{460}

$\alpha$ \deleted{\unsure{non} limite} on peut trouver un
\[
\deleted{\unsure{x}}\xi_\alpha\in E_{i_\alpha}\setminus\bigcup_{\beta<\alpha}E_{i_\beta}.
\]
L'application
\diagram{\begin{tikzcd}[ampersand replacement=\&,row sep=small,column sep=small]
\alpha\arrow[r,phantom,"\unsure{\mapstochar\relbar|\mapsto}" description]\&\deleted{\unsure{x}}\xi_\alpha\\
I_\alpha\arrow[r,phantom,"\unsure{-|\to}" description]\&E'_\infty
\end{tikzcd}}
est \deleted{\unsure{évi}}injective, donc $\operatorname{Card} E'_\infty\geq\operatorname{Card} I_{\alpha_0}$

\uline{Corollaire} 4.15.8.4. Pour \inserted{un $M$ et un $\pi_0$ donnés, pour} que le lemme-clef 4.15.8.1 soit valable, il f et s.
\noindent\makebox[0pt][r]{\smash{\rule[-10\baselineskip]{0.4pt}{10\baselineskip}}\quad}qu'il le soit \deleted{pour, \unsure{un}} dans le cas où $I=I_{\alpha_0}$, $\alpha_0$ étant le premier ordinal tel que $\operatorname{card}\alpha_0>\pi_0$.
Il revient au $=$ qu'on ne puisse trouver de système inductif $(\alpha_{ji})_{j\geq i}$ indexé par $I_{\alpha_0}$, à valeurs dans $M$, tel que pour \deleted{\unsure{j|j_0|j>}} $i<j$, on ait $\alpha_{ji}$ non inversible dans $M$.
% archive sheet number: 548
\scanpage{550}
\mspage{461}

Je vais résumer ce qui a été obtenu dans un énoncé récapitulatif :

\uline{Théorème} 4.15.9. Soient $C$ une \inserted{petite} catégorie, $\pi$ un cardinal, tel que $\pi\geq\operatorname{card} C$. Les conditions suivantes sont équivalentes (et je conjecture qu'elles sont toujours satisfaites) :

a) $\deleted{\unsure{C|M}}\unsure{\mathcal M|M}\overset{\mathrm{def}}{=}\mathrm{Ind}_\pi C$ est réduit à $\mathcal M_\pi$

a$'$) (Si $C$ est \unsure{Karoubienne|karoubienne}) Tout système inductif $\underline{x}=(x_i)_{i\in I}$ dans $C$, grand devant $\pi$, est essentiellement constant.

a$''$) Pour \unsure{tt} système inductif $(x_i)_{i\in I}$ dans $C$, grand devant $\pi$, il existe un foncteur cofinal $J\mathbin{\unsure{-|\to}}I$, avec $J$ \inserted{ens ordonné} \unsure{grd} devant $\pi$, et tel que le système réindexé \inserted{$(x_j)_{j\in J}$} satisfasse les
conditions suivantes : \unsure{1)|a)} tous les $x_j$ sont égaux à un même élément $x$ de $C$, \unsure{2)|b)} il existe un projecteur $p$ de $x$, tel que l'on ait $\alpha_{j'j}=p$ pour $j<j'$ $(j,j'\in J)$, où $\alpha_{j',j}:x_j\to\unsure{x_{j'}|x_j,|x_j}$ est la flèche de transition.

b) Pour tout $x\in\mathrm{Ob}\,C$, et pour tt monoïde quotient $M'$ de $M=\operatorname{End}_C(x)$, tel que $M'$ soit un monoïde de monomorphismes (i.e. $uv=uv'\Longrightarrow v=v'$ si $u,v,v'\in M'$), et pour tout système
inductif grand devant $\pi$ à valeurs dans $B_{M'}$, (défini par un système d'éléments $(\alpha_{ji})_{\substack{i\leq j\\i,j\in I}}$ d'un ens ordonné $I$ \unsure{grd} devant $\pi$, avec $\alpha_{hj}\alpha_{ji}=\alpha_{hi}$ pour $i\leq j\leq h$), il existe $i_0\in I$ tel
% archive sheet number: 549
\scanpage{551}
\mspage{462}

que $\alpha_{ii_0}=1$ pour $i\geq i_0$, i.e. $\alpha_{ji}=1$ pour $j\geq i\geq i_0$.

b$'$) $M'$ étant un monoïde de monomorphismes \unsure{comme dans|au sens de} b) (donc $\operatorname{card} M'\leq\operatorname{card} C\leq\pi$), \deleted{il n'existe pas} \inserted{soit de plus} $\alpha_0$ le plus petit ordinal \deleted{de cardinal} tel que $\operatorname{card}\alpha_0>\pi$, et soit $I=I_{\alpha_0}$ l'ens \inserted{ordonné} des ordinaux $\alpha<\alpha_0$ (lequel ens. est \unsure{grd} devant $\pi$).
\deleted{\unsure{J|I|Il}} Il n'existe \uline{pas} de système inductif dans $B_{M'}$ indexé par $I$, tel que \deleted{l'on ait $\alpha_{ji}$} pour $i<j$ dans $I$, $\alpha_{ji}$ soit un élément non inversible dans $M'$.
(\unsure{ou|donc} $u\mathbin{\unsure{\mapstochar\relbar|\mapsto}}\alpha_{ji}u:M'\mathbin{\unsure{-|\to}}M'$ est injectif mais non surjectif).

Signalons tout de suite le \deleted{\unsure{co}}

\uline{Corollaire} 4.15.10. Les conditions équivalentes précédentes sont satisfaites, si $\pi>\operatorname{card} C$.

En effet, soit $\alpha_1$ le \deleted{premier} plus petit ordinal tel que $\operatorname{card}\alpha_1>\operatorname{card} C$. \unsure{Donc|On} \deleted{\unsure{a}} $\pi_1\overset{\mathrm{def}}{=}\operatorname{card}\alpha_1$ \deleted{\unsure{est|$\xi$}} est le successeur de $\pi_0\overset{\mathrm{def}}{=}\operatorname{card} C$, et on a $\pi\geq\pi_1$, donc \deleted{$\operatorname{card}\alpha$} (si $\alpha_0$ est comme
explicité dans b$'$) ci-dessus) $\operatorname{card}\alpha_0>\pi_1$ donc $\alpha_1<\alpha_0$. Soit \deleted{\unsure{donc|des}} (avec les notations de b$'$), \deleted{$M_\alpha=$} Pour $\alpha\in I$ i.e. $\alpha<\alpha_0$, \unsure{l|1|d}
\[
M_\alpha=\underset{i<\alpha}{\varinjlim}\,\underset{\substack{\Vert\\M}}{M_i},
\]
on voit donc, \deleted{que} comme les $\alpha_{ji}$ sont injectifs et non surjectifs \deleted{\unsure{Card}}, que
\[
\operatorname{card} M_\alpha\geq\operatorname{Card}\alpha,
\]
% archive sheet number: 550
\scanpage{552}
\mspage{463}

et que pour $\alpha<\beta$, on a
\[
M_\alpha\hookrightarrow M_\beta\qquad\text{injectif}
\]
Appliquons ceci au cas où $\alpha=\alpha_1$, donc on trouve $\operatorname{card} M_{\alpha_1}\geq\operatorname{card}\alpha_1=\pi_1$ d'autre part en \unsure{prenant|posant} $\beta=\alpha_1+1$, d'où $M_\beta\simeq M$, on trouve une injection
\[
M_{\alpha_1}\hookrightarrow M
\]
d'où $\pi_1=\operatorname{card}M_{\alpha_1}\leq\operatorname{card}M=\pi_0$, une contradiction, ce qui établit b$'$.

\deleted{\uline{\unsure{Corollaire}}} \deleted{\uline{\unsure{Scholie|Lemme}}} \inserted{Corollaire} 4.15.11. Soit $C$ une catégorie essentiellement petite, $\pi$ un cardinal
tel que $\pi>\operatorname{card}\,\unsure{\mathrm{red}|\mathrm{réd}}\,C$. \deleted{Alors \unsure{est|M}} \inserted{Soit} $\unsure{\mathcal M|M}\overset{\mathrm{def}}{=}\mathrm{Ind}_\pi C$ \deleted{$\unsure{=\mathcal M_\pi|\mathcal M_\pi}$}, alors $\mathcal M=\mathcal M_\pi$, \deleted{et} c'est une enveloppe de Karoubi de $C$.

\unsure{OPS|OPJ} $C$ réduite, donc $\operatorname{card} C<\pi$, et on applique le \deleted{\unsure{th}} corollaire précédent.

Ainsi, le \deleted{\unsure{th}} résultat litigieux 4.8.13 (p. 212) est finalement prouvé sous une forme légèrement plus faible, avec $\pi>\operatorname{card}\,\mathrm{red}\,C$ \unsure{ald|al d} $\pi\geq\operatorname{card}\,\mathrm{red}\,C$. Du $=$ coup, les applications de 4.8.13 qu'on a faites jusqu'à présent (notamment le lemme 4.15.3).
% archive sheet number: 551
\scanpage{553}
\mspage{464}

4.16. \uline{Exemples et questions}

4.16.\deleted{\unsure{1}}0. J'avais crû que tout foncteur ``raisonnable'' (\unsure{ou|et} plus précisément, \unsure{tout} foncteur ``logiquement accessible'') entre catégories accessibles devait être accessible. Ce n'est malheureusement pas vrai, puisqu'un composé de deux foncteurs Hom
\[
\mathrm{Ens}\longrightarrow\mathrm{Ens}\qquad
E\to\underbrace{\mathrm{Hom}(\mathrm{Hom}(E,I),J)}_{J^{(I^E)}}
\]
est non accessible, si $I,J$ sont tous deux de cardinal $\geq2$. Prenant \unsure{p.ex.|p. ex.} $I=J=\{0,1\}$, on trouve le foncteur
\[
E\longmapsto\unsure{\mathscr P|\mathcal P|\mathfrak P}(\mathscr P(E))
\]
qui n'est pas accessible.

4.16.1. Deuxième désillusion : j'avais crû que pour \unsure{tte|tt} catégorie accessible $\unsure{\mathcal M|M}$, \unsure{toute|tte} sous-catégorie \inserted{$\unsure{\mathcal N|N}$} \unsure{strict\textsuperscript{t}|strictem\textsuperscript{t}} pleine et stable par facteurs directs,
``raisonnable'', \deleted{\unsure{devait}} (p.ex. ``logiquement accessible'') devait être accessible dans $\unsure{M|\mathcal M}$. Il n'en est rien, \deleted{\unsure{comme on|comme}} même dans le cas particulier de l'image inverse de la sous-catégorie accessible \unsure{des|de} $\mathrm{Ens}$
% archive sheet number: 552
\scanpage{554}
\mspage{465}

\[
\begin{aligned}
&\mathrm{Ens}^{*}\subset\mathrm{Ens}\\
&\quad\Vert\\[-1ex]
&\{\deleted{\unsure{x|A}}E\in\mathrm{Ens}\mid E\ne\varnothing\}
\end{aligned}
\]
par un contrafoncteur accessible \uline{représentable}
\diagram{\begin{tikzcd}[ampersand replacement=\&,column sep=large]
\unsure{\mathcal M|M}\arrow[r,mapsto]\&(\mathrm{Ens})^{\circ}\\
x\arrow[r,mapsto,no head]\&\mathrm{Hom}_{\mathcal M}(x,z),
\end{tikzcd}}
i.e la sous-catégorie \inserted{\deleted{\unsure{ouverte|ouvrante}} (\unsure{str. pl.}) $\mathcal M_z$} de $\mathcal M$ formée des objets $x$ tels que $\mathrm{Hom}_{\mathcal M}(x,z)\ne\varnothing$, n'est pas toujours accessible dans $\mathcal M$.
Ici, \deleted{\unsure{je}} ne suis pas parvenu à un exemple aussi élémentaire (cf \unsure{appendice|appendices}) -- j'ai
pris $\mathcal M=(k\text{-alg. comm.})$, $k$ un corps infini, \inserted{et $z$\unsure{,|{}} l'objet initial $k$ de $\mathcal M$} -- \deleted{mais l'exemple}
j'ai trouvé que $\mathcal M_z$ n'est pas \deleted{\unsure{stable}} $L_\pi$-stable dans $\mathcal M$, quelque soit le cardinal $\pi$ -- mais cet exemple utilisait
que $\mathcal U=\mathcal U_1$, \deleted{\unsure{premier}} \inserted{plus petit} univers contenant
$\unsure{\mathbb N|N}$.
\deleted{\unsure{avec $\mathbb N$|avec tt}} J'ignore si pour \unsure{tt|tout} autre univers,
on peut trouver un corps $k$ \inserted{de cardinal $\pi\in\mathcal U$} assez \unsure{grd|grand},
de façon que \unsure{chq|ch.} \deleted{\unsure{tout|tous}} $\pi$-ultrafiltre sur un ensemble $E\in\mathcal U$ soit trivial. Peut-être
% archive sheet number: 553
\scanpage{555}
\mspage{466}

un tel exemple, où \deleted{\unsure{la|le}} non la construction du couple $(\unsure{\mathcal M|M},\unsure{\mathcal N|N})$, mais la démonstration du fait que $\mathcal N$ n'est pas accessible dans $\mathcal M$, fait appel à des hypothèses particulières
sur l'univers ambiant, ne doit pas être considéré comme (tout à fait !) ``logiquement accessible''. Pour être vraiment probant, il faudrait un exemple de nature plus élémentaire.

Il reste la possibilité qu'il y ait un th. métamathématique qui assure que \unsure{tt|tout} \unsure{couple} $(\mathcal M,\mathcal N)$ logiquement accessible, avec
% slanted
\margin{C'est faux, \unsure{ex.|en|car} prendre une sous-catégorie fermée, cf \unsure{4.16.7.10|4.16.7 et 10|4.16.5.10|4.16.10} ci-dessous}
$\mathcal M$ accessible, et $\mathcal N$ sous-catégorie str. pl. $L_\pi$-stable dans $\mathcal M$ pour $\pi$ assez grand, soit accessible dans $\mathcal M$. Ceci s'appliquerait en particulier au cas d'une sous-catégorie \uline{\deleted{\unsure{ouverte}}\unsure{fermée|ouverte}} : toute sous-catégorie \deleted{fermée} \deleted{\unsure{de|str.}} (qui soit logiquement accessible) serait accessible dans $\mathcal M$.
% slanted
\margin{\deleted{\unsure{On|Cf} \illegible{2} sous-catégorie fermée, cf \unsure{4.16.7.1|4.16.7} \illegible{1}}}
L'exemple\inserted{-test} le plus simple serait celui du complémentaire \inserted{\unsure{fermé}} $\mathcal M'_z$ de la sous-catégorie ouverte $\mathcal M_z$ ci-dessus$^{*}$, i.e. la catégorie des
% slanted
\margin{NB elle est \unsure{acc.} dans $\mathcal M$ ssi $\mathcal M_z$ l'est !}
$x\in\mathrm{Ob}\,\mathcal M$ tels que $\mathrm{Hom}_{\mathcal M}(x,z)=\varnothing$, \unsure{ou|on} \unsure{non|encore|mor|un}
% archive sheet number: 554
\scanpage{556}
\mspage{467}

l'image inverse de la sous-catégorie fermée accessible \deleted{\unsure{de|déf}}
\[
e=\{\varnothing\}\subset(\mathrm{Ens})
\]
réduite à $\{\varnothing\}$, par le foncteur représentable
$x\unsure{\mapstochar\relbar|\mapsto}\mathrm{Hom}_{\unsure{\mathcal M|M}}(x,z):\mathcal M\to\mathrm{Ens}$. Il se pourrait \inserted{donc} qu'une telle catégorie soit \unsure{tjrs|tjrs.} accessible dans $\mathcal M$. Mais j'ai été incapable de le prouver, même dans le cas particulier ci-dessus où $\mathcal M=(k\text{-alg comm})$, $k$ étant un corps comm. infini. La difficulté est (dans tous les cas) la suivante : trouver,
pour $z$ fixé, un cardinal $\pi\inserted{>\pi_0}$ tel que pour \deleted{\unsure{tout}} \inserted{$x\in\mathrm{Ob}\,\mathcal M$, et tout} système (ou \uline{un} système) $\pi_0$-$\pi$-adapté $(x_i)_{i\in I}$ de limite inductive $x_{\infty}$ ($\pi_0$ un cardinal donné, utile pour $\mathcal M$),
\uline{si} $\mathrm{Hom}(x,z)$ \deleted{\unsure{est vide|est non vide}} $(\simeq\varprojlim\mathrm{Hom}_{\mathcal M}(x_i,z))$ est vide, alors $\exists\,i_0\in I$ tel que
$\mathrm{Hom}_{\mathcal M}(x_{i_0},z)\unsure{\ne\varnothing|=\varnothing}$
soit déjà vide. Dans l'exemple $\mathcal M=(k\text{-Alg comm})$, il faut donc prouver que \unsure{l'on|on} peut trouver un cardinal $\pi$ tel que pour \unsure{tte|toute} $k$-algèbre $A$ avec
$\mathrm{Hom}_{\unsure{\mathcal M|\mathcal M'}}(\deleted{\unsure{x}}A,k)=\varnothing$, il existe une sous-$k$-algèbre \inserted{$A'$} de cardinal $\leq\pi$, telle
% slanted
\margin{C'est faux, cf annotation marginale p. 466}
que $\mathrm{Hom}_{k}(\unsure{A'|A'_i},k)=\varnothing$. J'ai \deleted{\unsure{cru|eue}} tendance à croire que ça doit être vrai (et \unsure{=|en} \unsure{sans|sens}
% archive sheet number: 555
\scanpage{557}
% archive sheet number: 556
\mspage{468}

supposer que $k$ soit un corps, et aussi en prenant la condition $\mathrm{Hom}_{\unsure{\mathcal M|M}}(A,\unsure{Z|z})=\varnothing$, où $Z$ est un objet quelconque de $\mathcal M$. Ce genre de questions semble \unsure{étrangement|étonnamment} délicat\ldots

4.16.2. Il est possible, plus généralement, que la question suivante ait une réponse affirmative : Soit
\[
f:\mathcal M\unsure{\text{---}|\longrightarrow}\mathcal N^{\circ}
\]
un foncteur accessible, avec $\mathcal M,\unsure{\mathcal N|N}$ \inserted{cat.} accessibles, et soit $\underline T$ une sous catégorie de $\mathcal N$ accessible dans $\mathcal N$, et telle qu'il existe un cardinal $\pi$ ayant les propriétés
suivantes : $\mathcal N^\circ$ \deleted{\unsure{est}} satisfait $L_\pi$, \deleted{et $T^\circ$} (i.e. $\mathcal N$ stable par $\varprojlim$ grdes devant $\pi$)
et $T^\circ$ dans $\mathcal N^\circ$ est $L_\pi$-stable (i.e. $T$ stable dans $\mathcal N$ par $\varprojlim$ grdes \unsure{dev. à|dev.|devant} $\unsure{\pi|\pi_0}$).
Alors $\underline S\overset{\unsure{\mathrm{déf}|\mathrm{def}}}{=}f^{-1}(T)$ est-elle accessible dans $\mathcal M$ ?
% slanted
\margin{Faux, avec $\mathcal N=\mathrm{Ens}$ et un foncteur représentable $\mathcal M\to\mathrm{Ens}^{\circ}$, et prenant $\underline T=\{\varnothing\}$, cf \unsure{annotations marginales|annotation marginale} p. 466 et 467.}
Ceci généralise le cas $\mathcal N=(\mathrm{Ens})$, $\underline T=\{\varnothing\}$, $f$ foncteur représenté par $z\in\mathrm{Ob}\,\mathcal M$, envisagé \deleted{\unsure{tout|plus}} à l'heure. Avec le $=$ \deleted{\unsure{de|d}}type d'exemple, mais en prenant $\mathcal N=(\mathrm{Ens})^{\unsure{a|\alpha|*}}$
